\section{图像变换与函数复合}
\label{sec:02-Calculus/01-Functions-and-Precalculus/02}

特征处理脚本里有两行：一行取对数压缩长尾，一行做零均值标准化。

评审时有人问顺序能不能换。改一下试试：先标准化，一半样本变成负数；再取对数，整条流水线在第二步炸掉。反过来先取对数再标准化，一路畅通。同一个数据、同样两个函数，换个顺序就从可运行变成不可运行。

这不是脚本的毛病。把两个函数串起来，串出来的东西是一个新函数，而这个新函数的定义域由内外两层共同决定——外层能接受什么，取决于内层吐出什么。顺序换了，中间那批值就换了，能不能喂进下一步自然跟着变。

\subsection{复合是一条有类型的流水线}

写 \((g\circ f)(x)=g(f(x))\)，读作先做 \(f\) 再做 \(g\)，方向从右往左。要让这件事说得通，\(f\) 的输出必须是 \(g\) 认得的东西：\(f:A\to B\)、\(g:B\to C\) 时 \(g\circ f:A\to C\) 才成立。

于是 \((g\circ f)(x)\) 有定义需要两个条件同时满足，\(x\) 得在 \(f\) 的定义域里，而 \(f(x)\) 这个具体的值还得在 \(g\) 的定义域里。第一条谁都会查，第二条常被漏掉，因为它不看输入看中间值。

\begin{center}
\begin{tikzpicture}[x=1cm,y=1cm,font=\footnotesize,>=stealth]
  \draw[black!50] (0,0) rectangle (2.1,0.9);
  \node at (1.05,0.45) {输入 \(x\)};
  \draw[->] (2.25,0.45) -- (3.15,0.45);
  \node[above,font=\scriptsize] at (2.7,0.5) {\(f\)};
  \draw[black!55,fill=blue!10] (3.3,0) rectangle (5.9,0.9);
  \node at (4.6,0.45) {中间值 \(f(x)\)};
  \draw[->] (6.05,0.45) -- (6.95,0.45);
  \node[above,font=\scriptsize] at (6.5,0.5) {\(g\)};
  \draw[black!50] (7.1,0) rectangle (9.5,0.9);
  \node at (8.3,0.45) {输出 \(g(f(x))\)};
  \draw[black!55,dashed] (3.3,-0.15) -- (3.3,-0.85) -- (5.9,-0.85) -- (5.9,-0.15);
  \node[below] at (4.6,-0.9) {这一整块必须落在 \(g\) 的定义域里};
\end{tikzpicture}

\emph{定义域的约束不只挂在最左端；中间那个蓝块才是流水线最容易堵住的地方，而它的位置随 \(f\) 变化，换个顺序就换了一块。}
\end{center}

拿 \(h(x)=\sqrt{1-x^2}\) 验一遍。内层 \(f(x)=1-x^2\) 对所有实数都有输出，随便喂；外层 \(g(u)=\sqrt u\) 只收非负数。\(x=2\) 时内层给出 \(-3\)，这个值喂不进外层，于是 \(h\) 的实定义域是 \(\{x:1-x^2\ge0\}=[-1,1]\)。注意这个约束是从中间值那里读出来的，先化简再问定义域就会把它丢掉。三层以上的嵌套建议从最内层往外逐层写约束，漏掉外层限制是常见事故。

\subsection{平移和缩放不过是最简单的复合}

熟悉的抛物线 \(y=x^2\) 加工一下变成别的抛物线，这件事本身就是复合：先对输入做一次线性预处理，再交给核心函数，再对输出做一次线性后处理。写成一条模板就是

\[
g(x)=a\,f\bigl(b(x-h)\bigr)+k .
\]

输出端那一半符合直觉。\(+k\) 把整条曲线抬高 \(k\)，\(a\) 让它离横轴远 \(|a|\) 倍，\(a<0\) 再翻个个儿。你习惯从结果看问题，加就往上，乘就拉开。

输入端那一半会撞墙。第一反应是 \(f(x+2)\) 应该右移，毕竟加了个正数。但取 \(f(x)=x^2\)，想让新函数取到顶点值 \(0\)，需要括号里是 \(0\)，也就是 \(x=-2\)。加了 \(2\)，图像往左跑了 \(2\)。

\begin{center}
\begin{tikzpicture}[x=0.85cm,y=0.6cm,font=\small,>=stealth]
  \draw[->] (-4.7,0) -- (2.8,0) node[right] {\(x\)};
  \draw[->] (0,-0.8) -- (0,5.6) node[above] {\(y\)};
  \draw[blue!70!black,thick] plot[domain=-2.2:2.2,samples=60] (\x,{\x*\x});
  \draw[red!65!black,thick,dashed] plot[domain=-4.2:0.2,samples=60] (\x,{(\x+2)*(\x+2)});
  \fill (0,0) circle (2.2pt);
  \fill[red!65!black] (-2,0) circle (2.2pt);
  \draw[->,thick,black!75] (-0.15,-0.7) -- (-1.9,-0.7);
  \node[below,font=\footnotesize] at (-1.0,-1.1) {括号里加 \(2\)，顶点左移 \(2\)};
  \node[blue!70!black,font=\footnotesize] at (1.35,5.3) {\(y=f(x)\)};
  \node[red!65!black,font=\footnotesize] at (-3.0,5.3) {\(y=f(x+2)\)};
\end{tikzpicture}

\emph{横向变换作用在``喂给 \(f\) 的东西''上：要让 \(f\) 收到和原来一样的输入，新的 \(x\) 必须先补偿掉那次加法，所以符号看上去是反的。\(f(x-h)\) 才对应右移 \(h\)。}
\end{center}

缩放同理。\(f(2x)\) 要让内部取到原来的 \(u\)，需要 \(x=u/2\)，横坐标全部减半，图像被压扁；系数和效果互为倒数，这是横向变换的固定标志。真正该记住的不是``先横后纵''这类口诀，而是那句话：横向变换改的是输入，纵向变换改的是输出。遇到具体函数，挑几个特征点手算一遍变换后的坐标，比背顺序可靠。

顺带一提，这条模板只覆盖线性的预处理与后处理。内核本身被变形了（比如 \(f(x^3)\)）就套不上，得退回复合的一般框架。

\subsection{复合不可交换，这一点值钱}

打折与运费是最短的演示。设 \(f(p)=0.8p\)（打八折），\(g(q)=q+10\)（加运费）。先打折后加运费得 \(0.8p+10\)，先加运费后打折得 \(0.8p+8\)——运费也被打了折。\(p=100\) 时是 \(90\) 与 \(88\) 的区别，两套完全不同的商业规则。

在数据处理里这类差异更隐蔽也更贵。先剔除异常值再算均值，和先算均值再剔除异常值，得到的不是同一个统计量；训练集上先划分再标准化，和先标准化再划分，后者把验证集的信息泄漏进了训练。顺序不是风格问题，它是模型的一部分。

\subsection{会拆，是链式法则的前置技能}

复合的逆操作是分解。看到 \(h(x)=\sin\bigl((3x-1)^2\bigr)\)，要能立刻读出三层：最内 \(u=f(x)=3x-1\)，中间 \(v=g(u)=u^2\)，最外 \(r(v)=\sin v\)，于是 \(h=r\circ g\circ f\)。再看 \(k(x)=e^{\cos(x^2+1)}\)，同样是三层，多项式在最里，余弦居中，指数在最外。

现在不用求导。要练的是``谁包在谁外面''这个辨认动作，因为求导时局部变化率正是沿着这条链一层层往回传：

\[
h'(x)=r'(v)\cdot g'(u)\cdot f'(x).
\]

这条链在神经网络里就是反向传播，只不过每一层从一个数变成了一批数，中间量 \(u,v\) 变成了激活值，需要缓存下来供回传时使用（见 \hyperref[ch:124]{``损失与反向传播：目标怎样沿计算图返回参数''}）。前向是复合，反向是沿同一条复合链的逐层相乘，两件事共用一张分解图。

\subsection{对称性是写成等式的几何}

图像关于纵轴对称，写成 \(f(-x)=f(x)\)，称偶函数，\(x^2\)、\(\cos x\)、\(|x|\) 都是；图像关于原点对称，写成 \(f(-x)=-f(x)\)，称奇函数，\(x^3\) 和 \(\sin x\) 是。把几何性质变成代数等式的好处是它可以被验算，也可以省一半计算量：知道偶函数在 \([0,10]\) 上的行为，\([-10,0]\) 上的行为就不用再算。

两处容易出岔子。其一，讨论奇偶性的前提是定义域本身关于 \(0\) 对称，否则 \(f(-x)\) 根本没有定义。其二，绝大多数函数既不奇也不偶，\(x^2+x\) 就是，这很正常，不必硬往两类里塞。

回到本节开头那条流水线。我们已经会把函数串起来，也知道顺序和定义域怎样一起被牵动。剩下一个更硬的问题：串好之后能不能倒着走回去？把标准化撤销回原单位是可行的，把``取绝对值''撤销回原值就不行。下一节说明哪一类操作允许撤销，以及信息是在哪一步丢掉的。
